\documentclass[11pt]{article}


\setlength{\parindent}{0pt}
\usepackage{xltxtra}
\usepackage{hyperref}
\hypersetup{hidelinks}
\usepackage{url}
\urlstyle{tt}
\usepackage{xcolor}
\definecolor{CVBlue}{RGB}{23,110,191}
\usepackage{calc}
\usepackage{graphicx}
\usepackage{tikz}
\usetikzlibrary{calc}
\usepackage{fontspec}
\usepackage{xeCJK}
\usepackage{enumitem}
\CJKsetecglue{} %% 取消中文与数字之间的间隙


%% 主文档字体设置
\setmainfont[
    Path = fonts/Main/,
    Extension = .otf,
    BoldFont = texgyretermes-bold.otf, % 加粗字体
]{texgyretermes-regular.otf} % 正文字体

% 中文字体设置
\setCJKmainfont[
    Path = fonts/hansans/,
    Extension = .ttf,
    BoldFont = NotoSansSC-Bold.ttf, % 加粗字体
]{NotoSansSC-Regular.otf} % 正文字体


\usepackage{relsize}
\usepackage{xspace}

% 使用 fontawesome（部分图标）
\usepackage{fontawesome} 

% A4纸，上下左右边距
\usepackage[
    a4paper,
    left=1.2cm,
    right=1.2cm,
    top=1.5cm,
    bottom=1cm,
    nohead
]{geometry}

\renewcommand{\baselinestretch}{1.5} % 行间距设为1.5

\usepackage{titlesec}
\usepackage{enumitem}
\setlist{noitemsep} % 取消列表项间的额外间距
%\setlist{nosep} % 取消所有垂直间距
\setlist[itemize]{topsep=0.25em, leftmargin=*}
\setlist[enumerate]{topsep=0.25em, leftmargin=*}

% --- 用于控制【不同项目之间】的垂直距离 ---
\newlength{\interProjectSpacing}
\setlength{\interProjectSpacing}{0.9em} % <--- 在此调整项目之间的距离
\newcommand{\projectsep}{\vspace{\interProjectSpacing}}

% --- 用于控制【项目标题】与下方【项目描述】的距离 ---
\newlength{\intraProjectTitleSep}
\setlength{\intraProjectTitleSep}{0.4em} % <--- 在此调整标题和描述的距离
\newcommand{\titlebreak}{\\[\intraProjectTitleSep]}

% --- 用于控制【项目描述】与下方【要点列表】的距离 ---
\newlength{\intraProjectListTopSep}
\setlength{\intraProjectListTopSep}{0.2em} % <--- 在此调整描述和列表的距离

% =======================================================================


\titleformat{\section}         % 定制 \section 命令 
{\large\bfseries\raggedright} % 将 section 标题设置为大号、粗体且左对齐
{}{0em}                      % 可用于为所有 section 添加前缀（如“章节...”）
{}                           % 可用于在标题前插入代码
[{\color{CVBlue}\titlerule}]  % 在标题后插入一条横线
\titlespacing*{\section}{0cm}{*1.6}{*1.2}



\begin{document}
\pagenumbering{gobble}

%%%% 利用tikz来定位照片
%\begin{tikzpicture}[remember picture, overlay] 
%    \node[anchor = north east] at ($(current page.north east)+(-2cm,-1.2cm)$) {\includegraphics[height=3cm]{avatar.jpg}};
%  \end{tikzpicture}%
  %%%% 利用tikz来定位学校Logo，这里只在第一页显示
  \begin{tikzpicture}[remember picture, overlay] 
    \node[anchor = north west] at ($(current page.north west)+(0.5cm,+1.0cm)$) {\includegraphics[height=6cm]{zju.png}};
  \end{tikzpicture}%
\centerline{\LARGE\bfseries{吴学让}} 

\centerline{\normalsize{\faPhone\ 15858267209 \quad \faEnvelopeO\ \href{mailto:3230100823@zju.edu.cn}{3230100823@zju.edu.cn}}} 

%\centerline{\normalsize{\faGithubSquare\ \href{https://github.com/maksymilan}{https://github.com/maksymilan} \quad \faRssSquare\ \href{https://maksymilan.github.io/}{https://maksymilan.github.io}}} 
    
\section{\makebox[\widthof{\faGraduationCap}][c]{\color{CVBlue}\faGraduationCap}\ 教育背景}    
\textbf{浙江省杭州高级中学} \hfill 2020.9 -- 2023.9\\[0.3em] % 标题和正文间加一点距离
\textbf{浙江大学} \hfill 2023.9 -- 至今\\[0.3em] % 标题和正文间加一点距离
材料科学与工程（卓越人才培养班）\quad 大三\\
\textbf{主修均绩：}4.62 / 5.0\\
\textbf{累计排名：}4 / 103
\begin{itemize}[nosep]
    \item 核心专业课程：《材料物理》95 / 5.0、《材料表征1》95 / 5.0、《材料表征2》93 / 4.8《材料表征3》96 / 5.0、《材料性能1》96 / 5.0、《材料科学基础1》93/4.8、《材料工艺学1》96 / 5.0
\end{itemize}

%\section{\makebox[\widthof{\faUsers}][c]{\color{CVBlue}\faUsers}\ 项目经历}
\section{\makebox[\widthof{\faCogs}][c]{\color{CVBlue}\faCogs}\ 项目经历}
% --- 第一个项目 ---
% 将标题行末尾的 \\ 替换为 \titlebreak 命令
\textbf{1、国家级大学生创新训练项目\quad SAM预组装对刮涂法制备钙钛矿太阳能电池的影响机制研究}\quad 导师：薛晶晶教授 \hfill 2025.03 -- 2026.05 \titlebreak
项目简介：通过Py3-SiOx预组装复合体系\textbf{构建体相纳米异质界面}，\textbf{突破}传统自组装小分子的\textbf{有序堆积模式}，大幅增加化学界面以实现SAM层\textbf{牢固锚定}，有效改善纯SAM\textbf{刮涂工艺效率低下}的短板，为高效率、高稳定性钙钛矿太阳能电池的大面积、规模化生产提供重要借鉴。

% 在 itemize 的选项中，使用 topsep=\intraProjectListTopSep 来控制上边距
%\begin{itemize}[nosep, topsep=\intraProjectListTopSep]
%    \item \textbf{数据与后端}：运用 \textbf{Go} 语言及 \textbf{GORM} 框架，实现对全站“烂坑”帖的7x24小时全天候\textbf{抓取}与\textbf{入库}(MySQL)。
%    \item \textbf{可视化前端}：采用原生 \textbf{HTML/CSS/JS} 构建了“烂坑”实时监控\textbf{作战大屏}。实现了坑主成分分析、挖坑热力图、年度烂坑王榜单等\textbf{可视化图表}，并通过 \textbf{AJAX} 异步刷新，%确保数据洞察的实时性。
%    \item \textbf{核心产出}：全面掌握了从\textbf{数据采集、清洗、LLM标注到可视化分析}的“反烂坑”全链路技术。
%\end{itemize}

% 使用 \projectsep 命令来分隔两个项目
\projectsep

% --- 第二个项目 ---
\textbf{2、材料学院《创新实践》课程项目\quad SAM团簇抑制机制及其对钙钛矿太阳能电池性能的调控}\quad 薛晶晶教授 \hfill 2025.05 -- 至今 \titlebreak
项目简介：通过\textbf{调控}Py3-SiOx预组装复合体系中SiOx纳米颗粒的\textbf{粒径和表面配体}，实现SAM团簇的\textbf{有效抑制}，构建更为均匀致密的空穴传输层，从而提升钙钛矿太阳能电池的综合光电性能。
%\begin{itemize}[nosep, topsep=\intraProjectListTopSep]
%    \item \textbf{并发与调度}：后端采用 \textbf{Go} 语言，充分利用 \textbf{Goroutine} 和 \textbf{Channel} 的高并发模型，模拟\textbf{用户并发在线}，执行定时挖坑，确保挖坑行动的\textbf{隐蔽性}与\textbf{高效性}。
%    \item \textbf{实时监控与响应}：运用 \textbf{WebSocket} 协议实现了对目标帖子的\textbf{实时监听与交互}。确保在烂坑被识破的第一时间，Agent能光速完成\textbf{“lktp”}或\textbf{反向钓鱼}操作，展现出极高的AI博弈能力。
%    \item \textbf{可视化GUI}：前端使用 \textbf{React} 框架构建了\textbf{后台管理的可视化面板 (Admin Panel)}。可一键下达挖坑指令、动态调整Agent的\textbf{指令}、并实时监控数据。
%\end{itemize}

%\section{\makebox[\widthof{\faCogs}][c]{\color{CVBlue}\faCogs}\ 技术栈}
%\begin{itemize}[nosep]
%    \item \textbf{编程语言：} \textbf{Go}, Python, C++
%    \item \textbf{开发工具：} SSH, Git, Vim, MakeFile,LaTex
%    \item \textbf{操作系统：} Linux
%\end{itemize}

\section{\makebox[\widthof{\faGraduationCap}][c]{\color{CVBlue}\faList}\ 获奖情况}
\begin{itemize}
    \item 第十五届全国大学生数学竞赛（非数学A类）三等奖 \hfill 2023.12
    \item 校级优秀团员 \hfill 2024.05
    \item 浙江大学云峰学园“优秀大学生”荣誉称号 \hfill 2024.06
    \item 浙江省政府奖学金、浙江大学二等奖学金、材料学院丁子上奖学金 \hfill 2024.11
    \item 浙江省大学生物理（理论）创新竞赛三等奖 \hfill 2024.12
    \item 校级优秀团干部 \hfill 2025.06
    \item 阙端麟奖学金、浙江大学二等奖学金 \hfill 2025.11
\end{itemize}
    
\section{\makebox[\widthof{\faUsers}][c]{\color{CVBlue}\faInfo}\ 其他}
\begin{itemize}[parsep=0.5ex]
    \item \textbf{外语水平：} CET-4  569, CET-6  539
    \item \textbf{编程水平：}《C语言程序设计》93 / 4.8
    \item \textbf{专业软件：} 掌握 Origin、MATLAB、ImageJ 和 Materials Studio 等软件的基本操作，可开展数据分析、图像处理及材料模拟相关工作。
\end{itemize}
\end{document}
